\section{Verification record: chained immediate transfer}
\label{dist:app:verification}

Every quantitative claim of \cref{dist:sec:chained} and
\cref{dist:app:chained-transforms} was checked by exact simulation of the
chained process. Sampling is direct, because the event-type decomposition of
\cref{dist:sec:key-observation} makes the sequence of event types a sequence
of independent coin flips and the holding times exponential of known rate.
This appendix records the suite, its tally, and how to reproduce it.

The suite \texttt{verify\_chained\_transfer.py} samples the chained process
exactly and checks every quantitative claim across three parameter sets,
$(\lambda,\delta)\in\{(1,1),\,(1,0.1),\,(2,0.5)\}$: 28 checks, all passing.

\begin{center}
\renewcommand{\arraystretch}{1.05}
\footnotesize
\begin{tabular}{@{}clp{0.44\textwidth}c@{}}
\toprule
ID & Name & What is checked & Result \\
\midrule
A & Coefficient identity &
  $\binom{n+k-2}{k-1}$ obeys the recurrence (exact integers) & 1/1 \\
B & Rupture-size laws &
  $\pr{r_k=n}$ vs $10^6$ chains, $k=1..4$, SE-unit comparison & 3/3 \\
C & Moments & $\E{r_k}=1+k\lambda/\delta$, $\var{r_k}=k\rho/s^2$ & 3/3 \\
D & PGF & $\E{z^{r_k}}=z(s/(1-\rho z))^k$ (SE units) & 3/3 \\
E & First rupture time & histogram vs $\delta J(t)$ & 3/3 \\
F & Interval transform &
  fixed founders $k\in\{1,2,3,5\}$: empirical Laplace transform vs the
  series and its $\Fhyp$ form & 3/3 \\
F2 & Interval mixture &
  chain intervals $T_2,T_3$ (random founders) vs pmf-weighted $\Lap{T}$ & 3/3 \\
G & Second rupture time & empirical $\E{\ee^{-ut_2}}$ vs $\Lap{t_2}(u)$ & 3/3 \\
H/H2 & $t_2$ density and mean &
  exact hypoexponential-mixture density vs histogram
  (relative $L^2\approx3\%$); exact $\E{t_2}$ & 3/3 \\
I & Interval means &
  strictly decreasing; simulation vs the exact series & 3/3 \\
\midrule
 & \textbf{Total} & & \textbf{28/28} \\
\bottomrule
\end{tabular}
\end{center}

The $\Fhyp$ cross-check in test F activates only when \texttt{scipy} is
available; the series itself, which the simulation is compared against, needs
only \texttt{numpy}.

\paragraph{Reproduction.} The suite, its report, its metrics and its figures
are shipped with this chapter's source.
\begin{verbatim}
cd verification
python3 verify_chained_transfer.py
# writes verify_chained_transfer_report.txt, metrics, figures
\end{verbatim}

\paragraph{What is not covered.} The suite is confined to $\mu=0$, where the
event-type decomposition holds. The joint laws of $(t_k,r_k)$ for general
$\mu>0$, and the joint law of $(t_M,r_M)$ for the full $M$-cell chain, are
neither derived nor checked here; see \cref{dist:sec:open-problems}. For
$\mu=0$ the decomposition \cref{dist:eq:rkdecomp} extends the size laws to
all $M$ at once, since nothing in it depends on where a cell sits in the
chain.
